%% Guide to the Algorithmic Redistricting Portfolio
%% Track-based navigation guide for the 120+ paper research program

\documentclass[12pt]{article}

\usepackage{booktabs}
\usepackage{longtable}
\usepackage{hyperref}
\usepackage{geometry}
\usepackage[utf8]{inputenc}
\usepackage{times}
\usepackage{amsmath}
\usepackage{enumitem}
\usepackage{array}

\geometry{margin=1in}

\hypersetup{
    colorlinks=true,
    linkcolor=blue,
    urlcolor=blue
}

\title{\textbf{A Guide to the Algorithmic Redistricting Portfolio}}
\author{Giovanni Della Libera}
\date{2026}

\begin{document}

\maketitle
\tableofcontents

\newpage

% =============================================================================
\section{Overview}
% =============================================================================

Congressional redistricting in the United States suffers a legitimacy crisis:
partisan mapmakers manipulate district boundaries to predetermine election
outcomes, eroding public confidence in representative democracy. This portfolio
demonstrates that \textbf{recursive graph bisection}---a computational method
used for supercomputer workload distribution---provides an algorithmic
alternative that \emph{cannot be instructed to gerrymander because it cannot
access partisan data}.

Applied to all 50 states across three census decades (2000/2010/2020), the
algorithm produces 435 congressional districts with near-perfect population
equality, geographic contiguity, and geometric compactness. The research
program comprises \textbf{190+ papers across twenty active tracks (A--G,
I--U)}, spanning algorithm foundations (Track B), validation (Track C),
legal compliance (Track D), experimental extensions (Tracks E--F),
ensemble-based plan evaluation (Track G), incumbency and apportionment
(Tracks I--J), compactness and partisan fairness metrics (Tracks K--L),
community character weights (Track M), plan construction (Track T), and
search and optimisation (Track U), with Track A providing cross-portfolio
synthesis.

Papers in this program are at varying stages of completion. Track B
foundations, several Track G ensemble papers, the former Track H papers now
merged into Track U, and selected Track T and U papers are accepted or ready
for submission. Remaining papers across all tracks are in draft. The overall
research program is ongoing.

\subsection*{Four Core Contributions}

\begin{itemize}[leftmargin=*,noitemsep,topsep=4pt]
    \item \textbf{Philosophical}: \emph{Impossibility defense}---the algorithm
        cannot see partisan data, therefore cannot be instructed to gerrymander
        by construction. Compact neutral maps can still reflect geographic voter
        concentration, as shown in Track C.5.

    \item \textbf{Technical}: Edge-weighted graph partitioning achieves
        $+22$--$44\%$ compactness improvement over enacted maps, with seed
        variance below 2\% for 96\% of states.

    \item \textbf{Empirical}: National-scale validation across 50 states
        $\times$ 3 census decades (150 redistricting scenarios), with
        adversarial robustness testing across spatial resolution, statistical
        methodology, and ensemble comparison.

    \item \textbf{Legal}: The algorithm produces 137 majority-minority
        districts vs.\ 68 enacted (+69 surplus), with a 42\% geographic
        concentration threshold identifying the feasibility boundary for VRA
        compliance.
\end{itemize}


% =============================================================================
\section{Portfolio Architecture}
% =============================================================================

The portfolio is organized into twenty active tracks. Track A synthesizes all
other tracks; Tracks B--G generate the empirical, algorithmic, legal, and
ensemble results that the original synthesis draws on; Tracks I--M extend the
program to incumbency, apportionment, compactness, partisan fairness, and
community character weights; Tracks N--S extend to population counting,
democratic outcomes, reform pathways, 2030 planning, adversarial robustness,
and statistical inference; Tracks T--U factor the overloaded algorithm material
into plan construction and search/optimization.

\vspace{6pt}

\begin{longtable}{@{}llp{5.2cm}l@{}}
\toprule
\textbf{Track} & \textbf{Theme} & \textbf{Key Headline Result} &
    \textbf{Papers} \\
\midrule
\endhead

\textbf{A} & Portfolio Synthesis &
    190+ papers synthesized across 20 active tracks; 50 states $\times$ 3 decades
    $\times$ $+22\%$ compactness at CV $<$ 2\% &
    6 \\[4pt]

\textbf{B} & Algorithm Foundations &
    Recursive bisection, edge weighting, METIS formulation, complexity,
    n-way comparison, and seed-sensitivity foundations &
    9 \\[4pt]

\textbf{C} & Validation \& Robustness &
    Results robust across 130$\times$ unit-count range; near-zero efficiency
    gap as byproduct; 30\% more swing districts than enacted &
    10 \\[4pt]

\textbf{D} & VRA \& Legal Implementation &
    +69 majority-minority district surplus; 42\% geographic threshold;
    model federal statute formalizing ApportionRegions &
    9 \\[4pt]

\textbf{E} & Experimental Extensions &
    Geometric cost of federalism quantified; proportionality gain from
    multi-member variants; international applications &
    8 \\[4pt]

\textbf{F} & State Legislative Redistricting &
    Algorithm scales from 435 congressional to 400-district state chambers;
    NestSection achieves senate $= 2{\times}$ house spine compatibility &
    7 \\[4pt]

\textbf{G} & Ensemble Methods \& Comparison &
    Bisection plan at 0.1--0.7th compactness percentile (WI/GA/PA/CA),
    50th in NC; G.14 (practitioner ensemble comparison) accepted at 4.0/4 &
    16 \\[4pt]

\textbf{I} & Incumbency Effects &
    Incumbent pairing probability, safe-seat dynamics, and legal treatment
    of incumbency as redistricting criterion &
    5 \\[4pt]

\textbf{J} & Apportionment Methods &
    Huntington-Hill through Jefferson/D'Hondt comparison; bisect-apportion
    implementation verified against official 2020 apportionment &
    7 \\[4pt]

\textbf{K} & Compactness Taxonomy &
    Polsby-Popper, Reock, Convex Hull, Schwartzberg, Length-Width,
    population-weighted; multi-metric court guide &
    8 \\[4pt]

\textbf{L} & Partisan Fairness Metrics &
    Efficiency gap, mean-median, declination, seats-votes curve;
    proportionality paradox: sigma$\approx$0 for competitive states &
    7 \\[4pt]

\textbf{M} & Community Character Weights &
    Seven measurable signals (economic, land use, housing, commuting,
    topographic, administrative, transit) operationalizing communities
    of interest doctrine &
    9 \\[4pt]

\textbf{N} & Population Counting &
    Who counts in the balance metric: total pop vs.\ citizen VAP vs.\ prison-adjusted;
    N.5 definitive 50-state comparison; Evenwel v.\ Abbott implications &
    6 \\[4pt]

\textbf{O} & Outcomes \& Representation &
    Does compact redistricting improve democracy? RQI: bisect plans 0.62 vs.\ 0.41
    enacted; $+23$pp competitive districts; $-0.07$ DW-NOMINATE; $-14$min distance &
    6 \\[4pt]

\textbf{P} & Reform Pathways &
    Federal legislation (DIA), ballot initiatives (CA/CO/MI), commission integration,
    court-ordered algorithmic remedies, implementation costs ($<\$50$K) &
    6 \\[4pt]

\textbf{Q} & 2030 Forward Analysis &
    Demographic projections (TX$+3$, NY$-2$), data infrastructure (43 states need BG),
    TX $k=41$ prime fallback proved, ratio-optimal recommended &
    5 \\[4pt]

\textbf{R} & Adversarial Robustness &
    Five gaming vectors; combined max $< 0.5$ seats nationally; all detectable
    via SHA-256 audit chain; certificate pinning limitation documented &
    5 \\[4pt]

\textbf{S} & Statistical Inference &
    ESS-corrected permutation test (p=0.004 for NC at 10K steps); Bayesian Detection
    Score via Beta posterior; FDR control on 1{,}050 L-track tests &
    5 \\[4pt]

\textbf{T} & Plan Construction &
    GeoSection, AreaSection, ApportionRegions, VRASection, CVD, BFS Growth,
    and Moving-Knife structure methods factored out of the former B track &
    13 \\[4pt]

\textbf{U} & Search \& Optimization &
    ConvergenceSweep, simulated annealing, tempering, ILP certificates,
    Pareto/NSGA-II, and former Track H search strategies &
    11 \\

\bottomrule
\caption{Portfolio track summary (190+ papers across twenty active tracks A--G, I--U).}
\label{tab:tracks}
\end{longtable}

\subsection*{Three-Layer Compositor}

Every redistricting run is governed by three orthogonal choices that
correspond to the three-layer compositor of the \texttt{bisect} binary:

\begin{center}
\begin{tabular}{@{}lll@{}}
\toprule
\textbf{Flag} & \textbf{Controls} & \textbf{Key Variants} \\
\midrule
\texttt{--structure} & Bisection tree shape &
    \texttt{standard-bisect}, \texttt{prime-factor} (ApportionRegions),
    \texttt{ratio-optimal} (GeoSection), \texttt{nway} \\[4pt]
\texttt{--weights-override} & METIS edge/vertex signal &
    \texttt{geographic} (default), \texttt{county}, \texttt{unweighted},
    \texttt{vra-aligned} \\[4pt]
\texttt{--search} & Seed-selection strategy &
    \texttt{single}, \texttt{convergence} (T=600), \texttt{percentile},
    \texttt{bisection-ensemble}, \texttt{forest-recom}, \texttt{vra-recom} \\
\bottomrule
\end{tabular}
\end{center}

Track B addresses algorithm foundations; Track T addresses plan-construction
and structure variants; Track G addresses ensemble algorithms; Track U
addresses search, optimization, certification, and algorithm selection.


% =============================================================================
\section{Reading Paths}
% =============================================================================

The following curated reading paths guide each audience to the papers most
relevant to their interests. Papers are cited by track code (e.g., B.1, C.5).
Paper track codes correspond directly to the folder names under
\texttt{research/}.

% -----------------------------------------------------------------------------
\subsection{For Political Scientists}
% -----------------------------------------------------------------------------

\begin{enumerate}[leftmargin=*,noitemsep,topsep=2pt]
    \item \textbf{B.1} (Recursive Bisection): Establishes the impossibility
        defense---the algorithm cannot see partisan data---and demonstrates
        $+22\%$ compactness improvement over enacted maps across 50 states
        and three census decades.

    \item \textbf{D.0} (VRA Compliance): Documents the 137 vs.\ 68
        majority-minority district count, establishing that algorithmic
        redistricting exceeds current VRA performance by 69 districts without
        explicit racial targeting.

    \item \textbf{D.1} (Threshold Analysis): Identifies the 42\% minority
        concentration threshold as the geographic limit of algorithmic VRA
        compliance, with direct implications for the Gingles preconditions.

    \item \textbf{C.5} (Efficiency Gap Analysis): Shows that near-zero
        partisan efficiency gap emerges as a \emph{byproduct} of compactness
        optimization across all 50 states, without partisan input.

    \item \textbf{C.8} (Competitive Elections): Demonstrates that algorithmic
        maps produce 30\% more swing districts than enacted plans, addressing
        the competitive-elections objection directly.

    \item \textbf{T.4} (ApportionRegions): Presents the prime-factor
        bisection tree producing NC 7D/7R and 223D/209R nationally---the
        clearest demonstration that geographic structure drives partisan
        outcome.

    \item \textbf{A.0} (Synthesis Metapaper): Integrating account of all 60+
        empirical results, written for a \emph{Science}-level general audience.
\end{enumerate}

% -----------------------------------------------------------------------------
\subsection{For Computer Scientists}
% -----------------------------------------------------------------------------

\begin{enumerate}[leftmargin=*,noitemsep,topsep=2pt]
    \item \textbf{B.2} (Edge-Weighted Bisection): Introduces geographic
        proximity weights lifting compactness an additional $+22\%$ beyond
        B.1's baseline, explaining why single-objective edge-cut dominates.

    \item \textbf{B.3} (Multi vs.\ Edge): Proves that single-objective
        edge-cut dominates multi-constraint METIS for redistricting---the
        central algorithm-design finding of the program.

    \item \textbf{B.6} (Computational Complexity): Proves NP-hardness of
        balanced planar bisection and characterizes METIS runtime as
        $O(n_{\mathrm{iter}} \times n \log k)$ with empirical near-linear
        scaling ($O(n^{1.07})$).

    \item \textbf{B.7} (Solution Space): Establishes solution-space tightness:
        seed CV $<$ 2\% for 96\% of states, partisan variance near zero
        across all runs.

    \item \textbf{G.14} (Ensemble Comparison): Practitioner comparison of
        redistricting ensemble algorithms; highest-scored paper in the
        portfolio at 4.0/4.

    \item \textbf{U.10} (Redist-Ensemble): Presents the Rust ReCom
        implementation achieving 2500$\times$ speed over the Python
        GerryChain baseline.

    \item \textbf{T.10} (Centroidal Voronoi): Geometric district construction
        via CVD; accepted at 3.9/4. T.11 extends to Euclidean distance
        on EPSG:5070 projections for 2--6\% additional Polsby-Popper gain
        in coastal and elongated states.
\end{enumerate}

% -----------------------------------------------------------------------------
\subsection{For Legal Scholars}
% -----------------------------------------------------------------------------

\begin{enumerate}[leftmargin=*,noitemsep,topsep=2pt]
    \item \textbf{B.1} (Recursive Bisection): The impossibility defense and
        its implications under \emph{Rucho v.\ Common Cause}; Huntington-Hill
        as precedent for mathematical governance.

    \item \textbf{B.02} (One Federal Law): Formalizes ApportionRegions as the
        one-sentence geographic completion of Huntington-Hill; the model
        federal statute.

    \item \textbf{D.0} (VRA Compliance): Full VRA Section 2 compliance
        analysis: 137 vs.\ 68 majority-minority districts and the totality-of-
        circumstances framework.

    \item \textbf{D.1} (Threshold Analysis): The 42\% geographic threshold
        and its relationship to the Gingles compactness precondition.

    \item \textbf{D.3} (Compactness Tradeoff): VRA--compactness tradeoff
        characterized: real but bounded, with most VRA districts achievable
        within 85\% of baseline compactness.

    \item \textbf{D.4} (Legal Implementation): 50-state adoption pathway
        analysis; the model statute that formalizes ApportionRegions as the
        statutory vehicle.

    \item \textbf{D.5} (Gingles Bloc-Voting): Expert witness methodology for
        the Gingles bloc-voting analysis required when algorithmic plans face
        Section 2 challenge.
\end{enumerate}

% -----------------------------------------------------------------------------
\subsection{For Practitioners (Redistricting Commissions)}
% -----------------------------------------------------------------------------

\begin{enumerate}[leftmargin=*,noitemsep,topsep=2pt]
    \item \textbf{B.1} (Recursive Bisection): Core method overview;
        compactness and partisan neutrality at national scale.

    \item \textbf{B.2} (Edge-Weighted Bisection): Implementation details for
        compact districts using geographic proximity weights.

    \item \textbf{T.4} (ApportionRegions): Prime-factor bisection tree for
        balanced partisan outcomes; the algorithm most directly usable by
        commissions.

    \item \textbf{U.1} (Convergence Sweep): The statutory seed formula
        (T=600) balancing computational cost against partisan variance
        reduction.

    \item \textbf{U.8} (PercentileSweep): Statutory legal posture selection
        via percentile position in the ensemble---the practitioner-facing
        search strategy.

    \item \textbf{C.9} (Adoption Case Studies): Arizona, California, and
        North Carolina adoption pathway analysis; political feasibility
        documentation.

    \item \textbf{A.5} (Policy Brief): 4-page brief targeting legislative
        staff and redistricting commissioners; entry point for non-technical
        audiences.
\end{enumerate}

% -----------------------------------------------------------------------------
\subsection{For Judges and Legislative Staff}
% -----------------------------------------------------------------------------

\begin{enumerate}[leftmargin=*,noitemsep,topsep=2pt]
    \item \textbf{A.2} (Portfolio Summary): Rapid briefing document condensing
        the headline findings across all 8 tracks.

    \item \textbf{A.5} (Policy Brief): 4-page accessible summary for
        non-technical decision-makers.

    \item \textbf{B.02} (One Federal Law): The one-sentence geographic
        completion of Huntington-Hill and its constitutional grounding.

    \item \textbf{D.4} (Legal Implementation): Model federal statute and
        50-state adoption pathways.

    \item \textbf{D.5} (Gingles Bloc-Voting): The evidentiary methodology
        courts require when algorithmic plans face VRA challenge.

    \item \textbf{C.6} (User Study): Evidence that algorithmic maps are rated
        fairer by the public; legitimacy basis for adoption.
\end{enumerate}


% =============================================================================
\section{Glossary}
% =============================================================================

\begin{description}[leftmargin=2.2cm,style=nextline,noitemsep,topsep=4pt]
    \item[ApportionRegions] Prime-factor bisection tree structure that splits
        the state into sub-regions corresponding to the prime factors of the
        district count, enabling balanced geographic and partisan outcomes.
        Formalized as a one-sentence federal statute in B.02.

    \item[Bisection tree] The hierarchical sequence of binary splits that
        defines how a state is partitioned into districts under recursive
        bisection. Structure variants (standard, prime-factor, ratio-optimal)
        are the subject of Track T plan-construction sub-track.

    \item[Census tracts] Geographic units used as the atomic partitioning
        units ($\sim$4,000 residents each); the building blocks from which
        districts are assembled.

    \item[Compactness] A geometric measure of district shape regularity.
        Polsby-Popper ($4\pi \times \text{Area} / \text{Perimeter}^2$, range
        $[0,1]$) is the primary measure used throughout the portfolio;
        higher values indicate more compact districts.

    \item[Edge-weight] A numerical cost assigned to the graph edge connecting
        two adjacent census tracts. Higher weights make an edge expensive to
        cut, encouraging the partitioner to preserve adjacency between similar
        tracts, which produces more compact districts.

    \item[Efficiency gap] Partisan fairness measure: the difference in wasted
        votes between parties divided by total votes. A near-zero efficiency
        gap emerges as a byproduct of compactness optimization (C.5).

    \item[Ensemble] The distribution of all valid redistricting plans
        consistent with statutory constraints, sampled by MCMC methods
        (GerryChain, ReCom, forest-ReCom, etc.). Track G positions the
        bisection plan within this distribution.

    \item[GerryChain / ReCom] Open-source Python toolkit (Metric Geometry and
        Gerrymandering Group) for MCMC redistricting ensemble sampling.
        GerryChain implements ReCom (recombination) proposals. Track G
        compares bisection plans to GerryChain-generated ensembles.

    \item[Gingles test] Three-part legal test from \emph{Thornburg v.\ Gingles}
        (1986) establishing preconditions for a successful VRA Section 2 claim:
        (1) minority group sufficiently large and geographically compact;
        (2) minority group politically cohesive; (3) majority votes as a bloc
        to defeat minority candidates.

    \item[Impossibility defense] The core philosophical claim of the portfolio:
        the METIS recursive bisection algorithm cannot be instructed to
        gerrymander by construction because it is not given access to partisan
        data. It minimizes edge cuts in a geographic graph---a purely geometric
        objective. Compact neutral maps can still reflect geographic voter
        concentration, as shown in Track C.5.

    \item[Majority-minority district] A district in which minority voters
        constitute at least 50\% of the voting-age population (VAP), satisfying
        the Gingles compactness precondition for VRA Section 2 protection.

    \item[METIS] Multilevel $k$-way graph partitioning library (Karypis \&
        Kumar, 1998). The core computational engine used throughout the
        portfolio to bisect the census tract adjacency graph.

    \item[NestSection] Bisection tree structure enforcing senate $= 2\times$
        house spine compatibility for bicameral redistricting (T.6, F.2).

    \item[PercentileSweep] Search strategy (U.8) selecting among valid plans
        by their partisan percentile position in the ensemble, enabling
        commissions to choose a legal posture without targeting a specific
        partisan outcome.

    \item[Polsby-Popper] See \emph{Compactness}.

    \item[VRASection] Bisection structure variant (T.7) that aligns bisection
        cuts with minority geographic concentration, enabling VRA compliance
        post-\emph{Allen v.\ Milligan} (2023) and \emph{Callais} (2025).
\end{description}


% =============================================================================
\section{Code and Data}
% =============================================================================

\begin{itemize}[leftmargin=*,noitemsep,topsep=4pt]
    \item \textbf{Primary tool}: The \texttt{bisect} Rust binary, built from
        the \texttt{crates/} workspace in this repository. It is the
        production-grade implementation of all three-layer compositor options
        described in Section~\ref{sec:portfolio} and runs approximately
        213$\times$ faster than the archived Python pipeline.

    \item \textbf{Replication package}: See A.4 for the AEA-compliant
        replication materials package, including configuration files, census
        data download scripts, and output reproducibility documentation.

    \item \textbf{Census data}: U.S. Census Bureau redistricting files for
        2000, 2010, and 2020, downloaded via the \texttt{bisect fetch} command.
        Raw data is not committed to the repository ($\sim$55~GB).

    \item \textbf{METIS library}: \url{http://glaros.dtc.umn.edu/gkhome/metis/metis/overview}.
        The portfolio also includes a pure-Rust METIS implementation
        (\texttt{bisect-metis} crate) for portable binary distribution.

    \item \textbf{Contact}: \href{mailto:giodl@microsoft.com}{giodl@microsoft.com}
\end{itemize}

\end{document}
